%%%%%%%%%%%%%%%%%%%%%%%%%%%%%%%%%%%%%%%%%%%%%%%%%%

%% CV-XeLaTeX in English
%% Bastián González-Bustamante
%% Leiden University
%% https://github.com/bgonzalezbustamante/CV-XeLaTeX

%% Based on the following repositories:
%% Awesome CV LaTeX Template for CV/Resume
%% https://github.com/posquit0/Awesome-CV
%% Bastián González-Bustamante's CV-LaTeX Template
%% https://github.com/bgonzalezbustamante/CV-LaTeX
%% Carla Cisternas' CV-LuaLaTeX Template
%% https://github.com/carlacisternasg/CV-LuaLaTeX

%% LaTeX Project Public License v1.3c
%% https://github.com/bgonzalezbustamante/CV-XeLaTeX/blob/master/LICENSE.md

%%%%%%%%%%%%%%%%%%%%%%%%%%%%%%%%%%%%%%%%%%%%%%%%%%

\vspace{3mm}
\cvsection{Summary}

\begin{cvparagraph}

\textcolor{black}{I am a Postdoctoral Researcher in Computational Social Science at Leiden University. I hold a DPhil (PhD) in Politics from the University of Oxford. My work bridges comparative politics, government and computational social science, with a particular interest in relational structures, political careers and political behaviour. I have applied social network analysis and graph-based methods to study political elites, legislative careers and co-sponsorship, including analyses of network structure and homogeneity in legislatures and applications of algorithms to political elite networks. I have also used network analysis to examine the intellectual structure of academic fields. More broadly, I build large-scale relational and text-as-data pipelines, deploy AI and machine learning models, and apply causal inference strategies to complex observational data. My recent peer-reviewed articles have appeared in {\itshape Nature}, {\itshape Social Science Computer Review}, {\itshape Public Opinion Quarterly}, {\itshape World Development}, {\itshape Artificial Intelligence and Law}, and {\itshape Digital Journalism}, among others.}

\textcolor{black}{My current research combines large-scale behavioural, relational and textual data to study political actors, institutions and patterns of interaction. Recent work analyses more than 21 million social-media posts to examine amplification, contestation and toxic interactions around protest news, combining machine learning with causal inference. I am also working on the Terrorist Group Adaptation \& Lessons for Counterterrorism (TERGAP) project, funded by the European Union, to build a new dataset of counterterrorism events and identify patterns of adaptation and unintended effects of counterterrorism. In parallel, I contribute to COST Action CA22150 on executive-bureaucratic careers and serve as a Research Leader at the Oxford Computational Political Science Group.}
\vspace{1mm}
\end{cvparagraph}
